\subsection{A blocked family of rectangular Specht commutants}
\label{sec:rectangular-specht-barrier}

The full centralizers above retain the unrestricted exponent. A different,
smaller family suggested by the two-cut Schur observables does not at
rectangular heights two, three and four.
We work over $\mathbb C$ and use base-two entropy $H$.
For fixed $q\ge2$, let $E_{q,n}$ be the Specht module of shape $(n^q)$
for $S_N$, where $N=qn$, in an invariant real orthonormal basis. Put
\[
 d_{q,n}=\dim E_{q,n},\qquad
 \mathcal C_{q,n}=\operatorname{End}_{S_N}(E_{q,n}\otimes E_{q,n}),
\]
with diagonal $S_N$ action. Let $Q_{q,n}$ minimize
$\operatorname{rank}M+\operatorname{rank}N'$ subject to
$I=M+\Gamma N'$ and $M,N'\in\mathcal C_{q,n}$, where $\Gamma$ is
ordinary partial transpose in the chosen basis.
This is a symmetric-group commutant on a rectangular Specht tensor
square, not the full $\mathrm{SL}_q$ centralizer on a qudit tensor power.

\begin{proposition}[Rectangular Specht barrier]
\label{prop:rectangular-specht-barrier}
The map $\Gamma$ preserves $\mathcal C_{q,n}$. Put $d=q^2$,
$t=\lfloor N/(2d)\rfloor$, and
\[
 h_{N,q}=H\left(\frac{N-(d-1)t}{N},
                    \frac tN,\ldots,\frac tN\right).
\]
For every $N=qn\ge2d$,
\begin{equation}\label{eq:rectangular-specht-finite}
 Q_{q,n}\ge
 \frac{2^{Nh_{N,q}}}{(N+d)^{d(d-1)/2+d}}.
\end{equation}
Consequently, for each fixed $q$,
\begin{equation}\label{eq:rectangular-specht-exponent}
 \liminf_{n\to\infty}\frac{\log Q_{q,n}}{\log d_{q,n}}
 \ge\frac{h_{q^2}}{\log q},\qquad
 h_d=H\left(\frac{d+1}{2d},\frac1{2d},\ldots,\frac1{2d}\right).
\end{equation}
The lower bound is strictly larger than $\sqrt2$ for $q=2,3,4$.
\end{proposition}

\begin{proof}
For any real orthogonal $R$, partial transpose commutes with
conjugation by $R\otimes R$. Thus it preserves the indicated
commutant and is a self-adjoint involution for the Hilbert--Schmidt
inner product.

We extend the determinant witness of
Proposition~\ref{prop:binary-raw-surplus} to every rectangular width.
For $Z\in(\mathbb C^q)^{\otimes4}$, let $D_0(Z),D_1(Z)$ be the
$q^2$ by $q^2$ determinants of its $12|34$ and $14|23$ flattenings.
Define the degree-$q$ invariant
\[
 L(Z)=\sum_{i^{(1)},\ldots,i^{(4)}\in[q]^q}
 \left(\prod_{a=1}^4
    \epsilon_{i^{(a)}_1\cdots i^{(a)}_q}\right)
 \prod_{\ell=1}^q
    Z_{i^{(1)}_\ell i^{(2)}_\ell i^{(3)}_\ell i^{(4)}_\ell},
\]
where $\epsilon$ is the alternating tensor with
$\epsilon_{1\cdots q}=1$. At
$W_q=E_{13}(q)E_{24}(q)$, both determinants are nonzero and
$L(W_q)=(q!)^2\ne0$.
Write $n=2qt+r$, $0\le r<2q$. Then
\[
 P=(D_0D_1)^tL^r
\]
has degree $N$ and transforms by $\det^n$ at each local
$\mathrm{GL}_q$. Its real coefficient vector $v$ therefore lies in
the four local Schur sectors $(n^q)$ and has nonzero pairing with
$W_q^{\otimes N}$. This bounded residue factor is why the argument
applies to every sufficiently large $n$.

By Schur--Weyl duality, as stated in
\cite[Section~2, p.~2]{Manivel2010}, and the invariant orthogonal forms, the
resulting invariant four-party space identifies with
$\mathcal C_{q,n}$. The projected crossing is the vectorized
identity, up to a nonzero normalization, and exchanging parties
$2$ and $4$ is $\Gamma$. In particular $\langle v,I\rangle\ne0$.
On the first cut, multiplication by $D_0^t$ twists both Schur
factors by the determinant character, shifting every partition by
$(t^d)$. Hence $v$ is supported only on cut types $\nu_d\ge t$.
The same statement holds for $\Gamma v$, using $D_1^t$.

Decompose
\[
 E_{q,n}\otimes E_{q,n}
   =\bigoplus_\nu [\nu]\otimes\mathbb C^{m_\nu}.
\]
Only partitions with at most $d=q^2$ rows occur, because this tensor
square embeds in $(\mathbb C^{q^2})^{\otimes N}$.
Each commuting operator has blocks $I_{[\nu]}\otimes X_\nu$;
a nonzero block has rank at least $f^\nu=\dim[\nu]$.
If both $M$ and $N'$ had zero blocks on every type $\nu_d\ge t$,
the equation $I=M+\Gamma N'$ would give
\[
 0\ne\langle v,I\rangle
   =\langle v,M\rangle+\langle\Gamma v,N'\rangle=0.
\]
Thus the rank sum is at least the smallest $f^\nu$ on those types.
No multiplicity-free assumption is needed.

The hook formula \cite[Section~1, pp.~1--2]{Stanley2004rectangular},
written by cancelling the missing hook lengths in each padded row, reads
\[
 f^\nu=
 \frac{N!\prod_{i<j}(\nu_i-\nu_j+j-i)}
      {\prod_{i=1}^d(\nu_i+d-i)!}
 \ge\frac{\binom{N}{\nu_1,\ldots,\nu_d}}
            {(N+d)^{d(d-1)/2}}
 \ge\frac{2^{NH(\nu/N)}}{(N+d)^{d(d-1)/2+d}}.
\]
The last inequality is the multinomial type bound; it also follows
by expanding into successive binomial coefficients and bounding
each most likely type below by the reciprocal number of types.
For $\nu_d\ge t$, entropy is minimized at
$((N-(d-1)t)/N,t/N,\ldots,t/N)$, proving
\eqref{eq:rectangular-specht-finite}.
For fixed $q$, the same hook estimate on $(n^q)$ gives
$\log d_{q,n}=N\log q-O_q(\log N)$, while
$h_{N,q}\to h_{q^2}$. This proves
\eqref{eq:rectangular-specht-exponent}.

For completeness the three strict comparisons can be certified
without decimal approximations. First $5^5<2^{12}$ gives
$h_4>3/2>\sqrt2$. Next
$18^{90}>10^{50}2^{207}$ gives $h_9>23/10$, whereas
$\sqrt2\log3<(10/7)(8/5)=16/7<23/10$ because $3^5<2^8$.
Finally $17^{80}<2^{327}$ and $5657^2>8\cdot2000^2$ give
\[
 h_{16}=5-\frac{17}{32}\log17
   >\frac{7241}{2560}>\frac{5657}{2000}>2\sqrt2.\qedhere
\]
\end{proof}

At height two, the relevant Kronecker square is in fact
multiplicity-free \cite[Theorem~1; proof in Section~3]{Manivel2010},
but the block-rank argument above does not require this stronger
fact. The barrier concerns the full indicated commutant family.
It does not apply to unrestricted matrices or to the exponent-equivalent
centralizers of the preceding propositions, and it supplies no
uniform lower estimate for the BV spectral regularization.
